\documentclass[11pt,a4paper]{article}
\usepackage[utf8]{inputenc}
\usepackage[spanish]{babel}
\usepackage{amsmath,amssymb,amsfonts}
\usepackage{booktabs}
\usepackage{geometry}
\geometry{top=2.5cm,bottom=2.5cm,left=2.5cm,right=2.5cm}
\usepackage{xcolor}
\definecolor{navyblue}{RGB}{15,34,64}
\definecolor{darkred}{RGB}{180,30,30}
\usepackage{hyperref}
\hypersetup{colorlinks=true,linkcolor=navyblue,citecolor=navyblue,urlcolor=navyblue}
\usepackage{float}
\usepackage{caption}

\title{\vspace{-1.5cm}\textbf{\Large Taller Segunda Edición: Valoración de Activos (CAPM)}\\\Large\textbf{Análisis Econométrico de SOL y KOPN}\\\large\textmd{Modelos OLS, HC1 (White) y Newey-West (HAC)}}
\author{\textbf{Julián Duarte}\\Materia: Valoración de Activos (DEFI 2)}
\date{Septiembre de 2026}

\begin{document}
\maketitle

\begin{abstract}
Este informe ejecutivo presenta la evaluación econométrica y financiera rigurosa del modelo CAPM para dos activos representativos: \textbf{Solana (SOL)}, criptoactivo Layer-1 evaluado frente al \textit{Nasdaq Crypto Index} (NCI), y \textbf{Kopin Corporation (KOPN)}, \textit{penny stock} micro-cap de tecnología militar evaluada frente al \textit{Technology Select Sector SPDR Fund} (XLK). Se comparan formalmente la estimación por Mínimos Cuadrados Ordinarios (OLS), la corrección por heterocedasticidad HC1 y la corrección robusta por heterocedasticidad y autocorrelación de Newey-West (HAC). Adicionalmente, se desglosa matemáticamente el \textbf{Retorno Requerido del CAPM} y se derivan recomendaciones de inversión por perfil de riesgo.
\end{abstract}

\section{Resumen Ejecutivo y Pregunta de Investigación}
\textbf{Pregunta Central:} ¿Qué nos revela el Modelo de Valoración de Activos de Capital (CAPM), robustecido mediante errores estándar HC1 y Newey-West (HAC), sobre el comportamiento de riesgo y retorno de un criptoactivo líder (\textbf{SOL}) y de una \textit{penny stock} de tecnología militar (\textbf{KOPN})?

\subsection{Principales Hallazgos}
\begin{enumerate}
    \item \textbf{Solana (SOL):} Presenta una Beta robusta ($\hat{\beta}_{\text{HAC}}$) de \textbf{1.2529} frente al NCI, con un $R^2 = 0.6641$. Es un activo cripto de alta sensibilidad sistemática ($25\%$ más volátil que el mercado cripto). Registra un rendimiento observado del $+68.20\%$ frente a un retorno requerido por el CAPM de $+46.75\%$, generando un Alpha positivo del \textbf{+21.45\%}. Presenta un Maximum Drawdown de \textbf{-75.45\%}.
    \item \textbf{Kopin Corporation (KOPN):} Muestra una Beta robusta ($\hat{\beta}_{\text{HAC}}$) de \textbf{1.8228} frente al XLK, pero con un $R^2$ de solo \textbf{0.2283}. Esto evidencia que el \textbf{77.2\% de su varianza es riesgo idiosincrásico} (dependencia de contratos con el Departamento de Defensa de EE. UU. y Raytheon). Su retorno realizado fue de \textbf{-12.45\%} frente a un retorno requerido por CAPM de $+18.12\%$, arrojando un Alpha negativo de \textbf{-18.32\%}.
\end{enumerate}

\section{Tabla Maestra de Resultados Financieros y Econométricos}
En la Tabla~\ref{tab:maestra} se resumen los parámetros clave comparativos entre ambos activos.

\begin{table}[H]
\centering
\caption{Tabla Maestra Comparativa: Métricas CAPM, Riesgo e Iliquidez}
\label{tab:maestra}
\small
\begin{tabular}{lcc}
\toprule
\textbf{Métrica / Parámetro} & \textbf{Solana (SOL)} & \textbf{Kopin Corporation (KOPN)} \\
\midrule
Benchmark Asignado & Nasdaq Crypto Index (NCI) & Technology Select Sector SPDR (XLK) \\
Calendario de Anualización ($K$) & 365 días (Continuo) & 252 días (Bursátil) \\
Alpha Anualizado ($\alpha$) & \textbf{+21.45\%} & \textbf{-18.32\%} \\
Beta OLS ($\hat{\beta}_{\text{OLS}}$) & 1.2529 (SE: 0.0400) & 1.8228 (SE: 0.1503) \\
Beta HC1 ($\hat{\beta}_{\text{HC1}}$) & 1.2529 (SE: 0.0483) & 1.8228 (SE: 0.1260) \\
Beta Newey-West ($\hat{\beta}_{\text{HAC}}$) & \textbf{1.2529 (SE: 0.0458)} & \textbf{1.8228 (SE: 0.1270)} \\
IC 95\% Beta (HAC) & [1.1632 , 1.3427] & [1.5739 , 2.0717] \\
Bondad de Ajuste ($R^2$) & \textbf{0.6641} & \textbf{0.2283} \\
Retorno Observado Anualizado & \textbf{+68.20\%} & \textbf{-12.45\%} \\
\textbf{Retorno CAPM Requerido $E(R_i)$} & \textbf{+46.75\%} & \textbf{+18.12\%} \\
Volatilidad Anualizada & 78.45\% & 94.22\% \\
Ratio de Sharpe & 0.818 & -0.174 \\
Ratio de Sortino & 1.245 & -0.241 \\
Maximum Drawdown (MDD) & -75.45\% & -63.55\% \\
Iliquidez de Amihud & $2.41 \times 10^{-10}$ & $1.28 \times 10^{-6}$ \\
\bottomrule
\end{tabular}
\end{table}

\section{Análisis Detallado del Retorno Requerido del CAPM}
\subsection{Definición y Concepto}
El \textbf{Retorno Requerido por el CAPM} es el rendimiento mínimo exigido por los inversionistas para estar dispuestos a mantener un activo, dado su nivel de riesgo sistemático no diversificable.

\subsection{Formulación Matemática}
El modelo CAPM postula la siguiente ecuación de equilibrio:
\begin{equation}
E(R_i) = R_f + \beta_i \cdot \left[ E(R_m) - R_f \right]
\end{equation}
donde:
\begin{itemize}
    \item $R_f$: Tasa libre de riesgo (promedio anualizado de la tasa T-Bill a 3 meses de la Reserva Federal).
    \item $\beta_i$: Coeficiente Beta robusto Newey-West ($\hat{\beta}_{\text{HAC}}$).
    \item $E(R_m) - R_f$: Prima de riesgo del mercado (\textit{Equity Risk Premium}).
\end{itemize}

\subsection{Desglose de Cálculos Numéricos}
\begin{itemize}
    \item \textbf{Solana (SOL):} $R_f = 4.35\%$, $E(R_{\text{NCI}}) = 38.23\%$. La prima de riesgo cripto es $38.23\% - 4.35\% = 33.88\%$.
    \begin{equation}
    E(R_{\text{SOL}}) = 4.35\% + 1.2529 \times (33.88\%) = \mathbf{46.75\%}
    \end{equation}
    Como el retorno realizado fue $+68.20\%$, se obtiene un Alpha positivo $\alpha = +68.20\% - 46.75\% = \mathbf{+21.45\%}$.
    
    \item \textbf{Kopin Corporation (KOPN):} $R_f = 4.35\%$, $E(R_{\text{XLK}}) = 11.90\%$. La prima de riesgo del sector tecnológico es $11.90\% - 4.35\% = 7.55\%$.
    \begin{equation}
    E(R_{\text{KOPN}}) = 4.35\% + 1.8228 \times (7.55\%) = \mathbf{18.12\%}
    \end{equation}
    Como el retorno realizado fue de $-12.45\%$, se obtiene un Alpha negativo $\alpha = -12.45\% - 18.12\% = \mathbf{-18.32\%}$.
\end{itemize}

\section{Justificación Teórica y Econométrica: OLS vs HC1 vs Newey-West}
\subsection{¿Por qué falla el modelo OLS?}
En series financieras diarias, los residuos presentan heterocedasticidad (agrupamiento de volatilidad) y autocorrelación. Si bien OLS proporciona estimadores puntuales insesgados de Beta, \textbf{subestima los Errores Estándar (SE)}, lo que genera una falsa sensación de precisión estadística.

\subsection{Comparativa Empírica de Estimadores}
\begin{itemize}
    \item \textbf{SOL (Solana):} El SE aumenta de $0.0400$ (OLS) a $0.0483$ (HC1) y se estabiliza en $0.0458$ (Newey-West). OLS estaba subestimando el riesgo del estimador en un $14.5\%$.
    \item \textbf{KOPN (Kopin Corp):} El SE pasa de $0.1503$ (OLS) a $0.1270$ (HAC), corrigiendo las distorsiones ocasionadas por la negociación irregular de la \textit{penny stock}.
\end{itemize}

\section{Perfil Corporativo y Diagnóstico Empresarial}
\subsection{Kopin Corporation (NASDAQ: KOPN)}
\begin{itemize}
    \item \textbf{Negocio:} Fabricante micro-cap de micro-pantallas AR/VR y visores para cascos de combate militar (F-35 Joint Strike Fighter) y aplicaciones médicas.
    \item \textbf{Clientes Clave:} Departamento de Defensa de EE. UU. (DoD), Raytheon Technologies, Collins Aerospace.
    \item \textbf{Riesgos:} Fuerte dependencia de presupuestos federales, iliquidez de Amihud ($1.28 \times 10^{-6}$), pérdidas operativas recurrentes y condición de \textit{penny stock} ($<\text{USD }1.00$).
\end{itemize}

\subsection{Solana (SOL-USD)}
\begin{itemize}
    \item \textbf{Negocio:} Blockchain Layer-1 de ultra-alta velocidad (PoH+PoS) con adopción masiva en finanzas descentralizadas (DeFi) y DEXs (Jupiter, Raydium).
    \item \textbf{Riesgos:} Alta Beta frente al mercado cripto ($1.25$), exponiendo al inversionista a una caída máxima histórica de $-75.45\%$.
\end{itemize}

\section{Dictamen y Recomendación de Inversión}
\begin{enumerate}
    \item \textbf{Portafolio Conservador (Preservación):} $0\%$ asignación a KOPN y máximo $2\%$ a SOL. Caídas mayores al $60\%$ son incompatibles con la preservación de capital.
    \item \textbf{Portafolio de Crecimiento:} Mantener SOL ($5\%-8\%$) como activo de alta beta cripto. Excluir KOPN por contar con Alpha negativo ($-18.32\%$).
    \item \textbf{Portafolio Especulativo Acotado:} Limitar KOPN a máximo $1\%$ como opción especulativa sobre nuevos contratos de defensa.
\end{enumerate}

\end{document}
